Es posible, pues, limitarse al caso en que $B$ es una $A$-álgebra de
tipo finito y, por tanto, noetheriana. Sean $\mathfrak q_i$
$(1\leqslant i\leqslant n)$ sus ideales primos mínimos. Como $B$ se
supone \emph{reducido}, se identifica con un \emph{subanillo} del producto $B_0$
de los $B/\mathfrak q_i$. Si $C_0$ es la clausura integral de $A$
en $B_0$, se tiene $C=B\cap C_0$. Si se pone
$B'_0=B_0\otimes_A A'$ y $C'_0=C_0\otimes_A A'$, la planitud da
$C'=B'\cap C'_0$ ($\mathbf{0}_{\mathrm I}$, 6.1.3); basta, pues, demostrar
que $C'_0$ es la clausura integral de $A'$ en $B'_0$. Pero $C_0$ es
el producto directo de los $C_i$, clausuras integrales de $A$ en los
$B_i=B/\mathfrak q_i$. Por tanto, $C'_0$ es el producto directo de
los $C'_i=C_i\otimes_A A'$, y basta demostrar que $C'_i$ es la
clausura integral de $A'$ en $B'_i=B_i\otimes_A A'$. La cuestión se reduce así al caso en que $B$ es un dominio entero y una $A$-álgebra
de tipo finito. Si $\mathfrak p$ es el núcleo del homomorfismo
$A\to B$, se tiene también
$B'=B\otimes_{A/\mathfrak p}(A'/\mathfrak pA')$. Como el morfismo
$\operatorname{Spec}(A'/\mathfrak pA')\to
\operatorname{Spec}(A/\mathfrak p)$ es normal, se puede sustituir
$A$ por $A/\mathfrak p$ y $A'$ por $A'/\mathfrak pA'$, y suponer
en consecuencia que $A\subset B$.

